\documentclass[11pt]{article}
\usepackage[margin=1in]{geometry}
\usepackage{amsmath,amssymb,amsthm}
\usepackage[hidelinks]{hyperref}

\newtheorem{theorem}{Theorem}[section]
\newtheorem{proposition}[theorem]{Proposition}
\newtheorem{corollary}[theorem]{Corollary}
\newtheorem{lemma}[theorem]{Lemma}
\theoremstyle{remark}
\newtheorem{remark}[theorem]{Remark}

\newcommand{\C}{\mathcal{C}}
\newcommand{\tu}{\tau}
\newcommand{\Om}{\Omega}

\title{Erd\H{o}s problem 647: reductions, density bounds,\\ and the prime-tuples barrier}
\author{Scott D. Hughes\\ \normalsize Independent researcher}
\date{}

\begin{document}
\maketitle

\begin{abstract}
Let $\tu$ be the divisor function. Erd\H{o}s problem 647 asks whether there is an
integer $n>24$ with $\tu(n-k)\le k+2$ for all $1\le k<n$; Erd\H{o}s expected the
answer to be no. Write $\C(x)$ for the number of such candidates $n\le x$. We give
a reduction of the candidate set, the modular part of which is formally verified in
Lean 4: every candidate $n>84$ satisfies $2520\mid n$ and, writing $n=2520N$, the
class of $N$ modulo $46189$ lies in an explicit set of $41$; and every candidate
$n>24$ lies in one of two explicit prime-chain families. The families give, by an
elementary Brun sieve, $|\C(x)|\ll x/(\log x)^{7}$; a separate growing-window
argument gives $|\C(x)|\le x\exp(-(\log\log x)^{2-o(1)})$. Both are proved in the
companion manuscripts. We then show why density methods stop here. The relaxation
that drives them does not collapse, and its optimal output $x^{1-o(1)}$ cannot give
finiteness; in the computed range, the genuine condition leaves only the known
solution $24$, and the deepest near-misses are prime chains. A uniform bound
excluding arbitrarily long prime-chain near-misses would imply finiteness; such a
bound appears to require input of prime-tuples type. The results are density-zero,
not a resolution of the problem.
\end{abstract}

\section{The problem}

For a positive integer $m$ let $\tu(m)$ be the number of its divisors. Erd\H{o}s
problem 647~\cite{ErdosProblems647Forum} asks:
\begin{quote}
is there an integer $n>24$ such that $\displaystyle\max_{m<n}\bigl(m+\tu(m)\bigr)\le n+2$?
\end{quote}
Equivalently, is there $n>24$ with $\tu(n-k)\le k+2$ for every $1\le k<n$. The bound
$n+2$ is best possible, and $n=24$ works; Erd\H{o}s offered a prize for an example
$n>24$ and wrote that it is ``extremely doubtful'' that infinitely many exist,
suggesting in fact that $\max_{m<n}(\tu(m)+m-n)\to\infty$. Call an integer $n>24$
satisfying the condition a \emph{candidate}, and let
\[
  \C(x)=\#\{\,n\le x: n\text{ is a candidate}\,\}.
\]

This note collects, in one place, what the reduction of the candidate set buys. Two
density bounds on $\C(x)$ are stated and attributed to the companion manuscripts
\cite{HughesChains,HughesWindow}; the modular reduction underlying them is formally
verified \cite{ProofChain}. The remaining sections then locate the obstruction to
proving finiteness. None of this resolves the problem: every bound below is
density-zero, and the step from density-zero to finiteness is traced to a
prime-tuples-type obstruction.

\section{The reduction}

\subsection*{Stage 1: modular constraints}

The divisor conditions at the shifts $k$ with small $\gcd(2520,k)$ force divisibility
and exclude residue classes. The following is checked over the explicit finite search
and is formally verified in Lean 4 against Mathlib; the verification artifact is the
release \texttt{v0.1.0} of \cite{ProofChain}.

\begin{theorem}[Stage-1 reduction; formally verified]\label{thm:stage1}
Every candidate $n>84$ satisfies $2520\mid n$. Writing $n=2520N$, the residue of $N$
modulo $M=11\cdot13\cdot17\cdot19=46189$ lies in an explicit set of $41$ classes.
\end{theorem}

The sieve modulo $11,13,17,19$ first leaves $96$ surviving classes; small-shift
divisor-budget arguments close $55$ of them, leaving the stated $41$. The full
deduction and its formalization are in \cite{ProofChain}; the open $41$-class set is
listed there.

\subsection*{Stage 2: prime-chain families}

\begin{theorem}[Prime-chain reduction]\label{thm:stage2}
Every candidate $n>24$ lies in one of the two families
\[
  n=8s+8,\quad s,\ 2s+1,\ 4s+3,\ 8s+7\ \text{all prime},
\]
or
\[
  n=16s+8,\quad s,\ 4s+1,\ 8s+3,\ 16s+7\ \text{all prime}.
\]
\end{theorem}

\begin{proof}[Sketch]
The conditions at $k=1,2$ force $n=2q+2$ with $q$ and $2q+1$ prime. Writing
$q-1=2^a m$ with $m$ odd, the condition at $k=4$ forces $q=2p+1$ with
$p,2p+1,4p+3$ prime. The condition at $k=8$, with $p-1=2^b s$, leaves exactly $b=1$
and $b=2$, the two displayed branches; the remaining cases are eliminated by the
finite divisor-budget check carried out in \cite{HughesChains}.
\end{proof}

\section{Density bounds}

Both families are admissible tuples of linear forms, so Brun's upper-bound sieve
applies. Once $2520\mid n$ (Theorem~\ref{thm:stage1}) the factor $3\mid n$ is
available, and the three further shifts $k\in\{3,6,12\}$ — whose $2520$-forced part is
exactly $k$ — adjoin three more prime forms to each family. The seven prime-forcing
shifts $\{1,2,3,4,6,8,12\}$ give, in each branch, an admissible seven-tuple
(root-union $1,2,4,6$ modulo $2,3,5,7$, and at most $7<p$ classes for $p>7$).

\begin{theorem}[Seven-form bound \cite{HughesChains}]\label{thm:seven}
$|\C(x)|\ll x/(\log x)^{7}$, unconditionally, with an effective absolute implied
constant.
\end{theorem}

The candidate condition can also be used only through the elementary implication
$\tu(N)\ge\Om(N)+1$, where $\Om$ counts prime factors with multiplicity. For a
window of length $H$ this gives $\Om\!\bigl(\prod_{k=1}^{H}(n-k)\bigr)\le H(H+3)/2$
for every candidate, and a Rankin moment together with the upper-bound
$\beta$-sieve fundamental lemma yields a bound that improves on every fixed power of
$\log x$. The companion growing-window manuscript proves the following.

\begin{theorem}[Growing-window bound \cite{HughesWindow}]\label{thm:t3}
$|\C(x)|\le x\exp\!\bigl(-(\log\log x)^{2-o(1)}\bigr)$.
\end{theorem}

\begin{remark}
Theorem~\ref{thm:t3} supersedes Theorem~\ref{thm:seven} and the earlier
$x/(\log x)^{k}$ bounds: $\exp(-(\log\log x)^{2-o(1)})$ is smaller than
$(\log x)^{-A}$ for every fixed $A$. We keep Theorem~\ref{thm:seven} because its
proof is elementary and self-contained.
\end{remark}

\section{The finiteness barrier}

Theorems~\ref{thm:seven} and~\ref{thm:t3} are density-zero results. We now explain
why no refinement of this kind reaches finiteness, and where the genuine difficulty
lies.

\subsection*{The density method cannot reach finiteness}

The growing-window bound of Theorem~\ref{thm:t3} uses the candidate condition only
through $\Om(n-k)\le k+1$. For a window length $H$ set
\[
  N_{\Om,H}(x)=\#\{\,n\le x: \Om(n-k)\le k+1\ \text{for }1\le k\le H\,\}\ \ge\ |\C(x)|.
\]
Even a bound as strong as
$x\exp(-c(\log\log x)^2)$, the natural frontier suggested by this relaxation, would
still not imply finiteness.

\begin{proposition}\label{prop:ceiling}
$x\exp(-c(\log\log x)^2)=x^{1-o(1)}\to\infty$ for every $c>0$. Hence no bound of this
form, however sharp, implies $\C(x)=O(1)$.
\end{proposition}

\begin{proof}
$(\log\log x)^2=o(\log x)$, so $c(\log\log x)^2/\log x\to0$ and the exponent of $x$
tends to $1$.
\end{proof}

Direct enumeration confirms that the relaxation itself does not collapse:
$N_{\Om,24}(10^7)=208828$, and the count shows no sign of thinning to a finite set,
whereas the genuine condition is far more rigid.

\subsection*{Enumeration of the genuine condition}

Let $D(n)$ be the largest $D$ with $\tu(n-k)\le k+2$ for all $1\le k\le D$ (the depth
of the longest run of satisfied conditions starting at $n$). A candidate is exactly
an $n>24$ with $D(n)\ge n-1$. Enumeration to $10^7$ gives the following, computed
independently with two divisor-function implementations.

\begin{center}
\begin{tabular}{rl}
\hline
$D$ & $\#\{24<n\le 10^7: D(n)\ge D\}$\\
\hline
$1$ & $665014$\\
$2$ & $30653$\\
$4$ & $159$\\
$6$ & $2$\\
$7$ & $1$\\
$8$ & $0$\\
\hline
\end{tabular}
\end{center}

No $n>24$ up to $10^7$ survives eight consecutive conditions. The record depths grow
far more slowly than $n$: the deepest runs below $10^7$ reach
\[
  D(n)=4,\,5,\,6,\,7\quad\text{at}\quad n=720,\ 58680,\ 7224840,\ 9339120,
\]
the integer attaining each new depth being far larger than the last. A candidate
needs $D(n)\ge n-1$, whereas the deepest run anywhere below $10^7$ has $D=7$. The gap
between the required depth and the observed record widens with $n$, the numerical
form of Erd\H{o}s's $\max_{m<n}(\tu(m)+m-n)\to\infty$.

\subsection*{What finiteness requires}

The deepest near-misses are prime chains. At each record above, $n-1$ is prime,
$n-2=2p$, $n-3=3p'$ and $n-4=2^2p''$ with $p,p',p''$ prime; deeper shifts admit other
low-divisor shapes, and the run ends when a forced quotient first carries too many
divisors. So large depth at these near-misses is governed by simultaneous prime-chain
conditions of the type in Theorem~\ref{thm:stage2}.

This places the two directions of the problem on closely related prime-pattern
objects. The existence of arbitrarily deep near-misses is naturally explained by
prime-tuples heuristics, and Schinzel-type hypotheses predict such finite-window
behaviour. Finiteness would require the opposite kind of input: a uniform obstruction
preventing these prime-chain patterns from extending far enough. Such an obstruction
is not supplied by the density estimates above, and appears to lie in the same
prime-tuples barrier discussed for related divisor problems by Tao and
Ter\"av\"ainen~\cite{TaoTeravainen2025}. The reduction here makes that barrier
visible; it does not cross it.

\section{Formal verification and reproducibility}

Theorem~\ref{thm:stage1} is formalized in Lean 4 against Mathlib in the public
repository \cite{ProofChain}, release \texttt{v0.1.0}. The build is free of
\texttt{sorry}; an axiom audit accompanies the release, recording the single
project-level axiom and the standard foundations on which the verified statements
depend. The enumeration of Section 4 and the admissibility checks of Section 3 are
reproduced by the scripts in the same repository.

\section*{Acknowledgments}

The partial modular observations behind Theorem~\ref{thm:stage1} are due to
S.~Dutta and B.~Alexeev in the erdosproblems.com discussion of the
problem~\cite{ErdosProblems647Forum}, and the clean shift $k=3$ underlying the
seven-form bound (Theorem~\ref{thm:seven}) is due to K.~Kitamura. The formal
verification uses Lean 4 and Mathlib; numerical and symbolic checks used PARI/GP,
Wolfram, and SymPy. AI-based tools assisted with drafting and exploratory checking,
and all mathematical content was verified independently.

\begin{thebibliography}{9}

\bibitem{ErdosProblems647Forum}
T.~F.~Bloom,
\emph{Erd\H{o}s problem \#647},
\url{https://www.erdosproblems.com/647}, accessed 2026-06-03.

\bibitem{HughesChains}
S. D. Hughes,
\emph{Erd\H{o}s problem 647 and prime chains: an unconditional density-zero bound},
submitted for publication, 2026.

\bibitem{HughesWindow}
S. D. Hughes,
\emph{An unconditional density bound for Erd\H{o}s problem 647},
submitted for publication, 2026.

\bibitem{ProofChain}
S. D. Hughes,
\emph{Erd\H{o}s 647 proof chain} (Lean 4 formalization), release \texttt{v0.1.0},
\url{https://github.com/scottdhughes/erdos647-proof-chain}.

\bibitem{TaoTeravainen2025}
T. Tao and J. Ter\"av\"ainen,
\emph{Quantitative correlations and some problems on prime factors of consecutive
integers}, arXiv:2512.01739 (2025).

\end{thebibliography}

\bigskip
\noindent
Scott D. Hughes\\
Independent researcher\\
Email: \href{mailto:v@deltagray.com}{\texttt{v@deltagray.com}}

\end{document}
